\section{Introduction}
\label{sec:introduction}

Machine learning systems must balance task quality with computational
efficiency. Improvements at the level of an individual operation do not
necessarily translate into faster end-to-end execution, while increasing
the computational resources available to a model does not necessarily
improve its outputs. Building on Transformer-based ASR \citep{vaswani2017attention,radford2023robust} and retrieval-augmented generation \citep{lewis2020retrieval}, this report investigates these trade-offs through
two complementary workloads: GPU-accelerated automatic speech
recognition (ASR) and retrieval-augmented generation (RAG).

For ASR, we implement and evaluate CUDA operations used in the supplied
inference pipeline, including rotary positional embeddings, reductions,
matrix multiplication, and attention. We compare custom implementations
against PyTorch-backed alternatives and investigate the effects of
matrix-multiplication tile size and kernel selection. Our central
question is whether custom CUDA kernels improve end-to-end inference
latency while preserving transcription correctness.

For RAG, we investigate retrieval depth, evidence selection,
and the relationship between context size, answer quality, and
inference cost. We evaluate fixed context sizes, adaptive
similarity-gap policies, and an oracle evidence condition on
125 development queries. We measure retrieval effectiveness,
automated answer quality, prompt-token usage, and generation
latency.

The fully custom CUDA implementation was approximately 4.38 times
slower than the PyTorch baseline. A hybrid configuration using
PyTorch GEMM and attention achieved 1296.2 ms, compared with
1339.2 ms for PyTorch-only execution on the tested workload.
This small measured difference does not establish a general
performance advantage. Profiling and controlled substitutions
identified custom GEMM and attention as the principal bottlenecks.

In the 125-query RAG study, judged answer quality was broadly
similar across fixed context sizes of 3 to 20 documents.
Increasing context from 5 to 20 documents raised prompt-token
usage by approximately 4.1 times without an established quality
improvement. The higher oracle-evidence score suggests that
evidence availability is an important limitation, although
automated judging and single-generation evaluation constrain
the conclusions.

Together, the experiments illustrate a common systems principle:
greater implementation complexity or computational expenditure does
not guarantee better end-to-end outcomes. Effective optimisation
requires workload-specific measurement, controlled comparisons, and
explicit consideration of both quality and cost.
